\documentclass[11pt,a4paper]{article}

% ---- packages ----
\usepackage[T1]{fontenc}
\usepackage[utf8]{inputenc}
\usepackage{newpxtext,newpxmath}
\usepackage[margin=2.3cm]{geometry}
\usepackage{xcolor}
\usepackage{booktabs}
\usepackage{tabularx}
\usepackage{enumitem}
\usepackage{titlesec}
\usepackage{fancyhdr}
\usepackage{fontawesome5}
\usepackage[most]{tcolorbox}
\usepackage{needspace}
\usepackage{hyperref}

% ---- colours (match slides/custom.scss) ----
\definecolor{lsenavy}{HTML}{172C83}
\definecolor{lightnavy}{HTML}{EEF0F8}
\definecolor{midgrey}{HTML}{5A5A5A}

\hypersetup{colorlinks=true, linkcolor=lsenavy, urlcolor=lsenavy, citecolor=lsenavy}

% ---- headings ----
\titleformat{\section}{\Large\bfseries\color{lsenavy}}{}{0em}{}[\vspace{-0.4em}{\color{lsenavy}\rule{\textwidth}{0.6pt}}]
\titleformat{\subsection}{\large\bfseries\color{lsenavy}}{}{0em}{}
\titlespacing*{\section}{0pt}{1.4em}{0.8em}
\pretocmd{\section}{\needspace{10\baselineskip}}{}{}
\titlespacing*{\subsection}{0pt}{1.0em}{0.4em}

% ---- lists ----
\setlist{itemsep=0.15em, topsep=0.3em}
\setlist[itemize,1]{label={\color{lsenavy}\small$\bullet$}}

% ---- header / footer ----
\pagestyle{fancy}
\fancyhf{}
\renewcommand{\headrulewidth}{0pt}
\fancyfoot[L]{\small\color{midgrey}MY580 $\cdot$ Causal Inference in the Social Sciences}
\fancyfoot[R]{\small\color{midgrey}\thepage}

% ---- helpers ----
\newcommand{\blockmeta}[2]{\noindent{\color{midgrey}\faClock[regular]\ #1 \quad \faChalkboardTeacher\ #2}\par\smallskip}
\newtcolorbox{reading}{colback=lightnavy, colframe=lightnavy, arc=2pt, boxrule=0pt,
  left=8pt, right=8pt, top=5pt, bottom=5pt, before skip=0.6em, after skip=0.6em}
\newtcolorbox{labbox}[1]{colback=white, colframe=lsenavy, arc=2pt, boxrule=0.6pt,
  left=8pt, right=8pt, top=5pt, bottom=5pt, before skip=0.6em, after skip=0.6em,
  title={\faLaptopCode\ \ #1}, coltitle=white, colbacktitle=lsenavy, fonttitle=\bfseries}
\newcommand{\TheEffect}{\emph{The Effect}}

\begin{document}

% ================= TITLE =================
\begin{center}
  {\color{lsenavy}\rule{\textwidth}{1.2pt}}\\[0.8em]
  {\LARGE\bfseries\color{lsenavy} Causal Inference in the Social Sciences}\\[0.3em]
  {\large A Design-First Introduction}\\[0.6em]
  {\color{midgrey} MY580 Workshop $\cdot$ London School of Economics}\\[0.8em]
  {\color{lsenavy}\rule{\textwidth}{1.2pt}}
\end{center}

\vspace{0.6em}
\noindent
\begin{tabularx}{\textwidth}{@{}l X@{}}
  \textbf{Instructor} & Zach Dickson, Fellow in Quantitative Methodology \\
  \textbf{Format}     & One day: 4 contact hours plus a 1-hour lunch; lectures with two hands-on labs in R \\
  \textbf{Core text}  & Huntington-Klein, N. \emph{The Effect: An Introduction to Research Design and Causality}. Chapman \& Hall/CRC. Free online at \url{https://theeffectbook.net} \\
  \textbf{Materials}  & \url{https://github.com/z-dickson/my580-causal-inference} \\
  \textbf{Audience}   & MSc and PhD students who know basic regression. No prior causal inference needed. \\
\end{tabularx}

% ================= OVERVIEW =================
\section{Overview}

Most interesting questions in the social sciences are causal. Does canvassing increase turnout?
Do elite cues change citizens' behaviour? Did a policy reduce crime? This workshop is a gentle introduction
to \emph{thinking causally}. The emphasis is on research design rather than statistical machinery: what
we would need to observe to answer a causal question, why the ideal comparison is usually
impossible, and how experiments and clever observational designs try to get close to it.

The day moves from foundations (potential outcomes and the fundamental problem of causal inference),
to randomized experiments, to observational studies, and ends with a short tour of more advanced
topics. Each of the two main blocks ends with a short lab where students run and interpret code.

\subsection{Learning outcomes}
By the end of the workshop, students should be able to:
\begin{enumerate}[leftmargin=1.6em]
  \item Turn a social science question into a causal question with a clear counterfactual and estimand.
  \item Explain the fundamental problem of causal inference, and why a naive comparison mixes up the causal effect with selection bias.
  \item Explain \emph{why} randomization identifies causal effects, and spot common threats to experiments.
  \item Use directed acyclic graphs (DAGs) to state their assumptions and to find confounders and colliders.
  \item Explain the identifying assumption behind difference-in-differences, regression discontinuity and instrumental variables, and run and interpret a simple difference-in-differences analysis.
\end{enumerate}

% ================= SCHEDULE =================
\clearpage
\section{Schedule at a glance}

\noindent
\renewcommand{\arraystretch}{1.25}
\begin{tabularx}{\textwidth}{@{}l X l@{}}
  \toprule
  \textbf{Time} & \textbf{Block} & \textbf{Format} \\
  \midrule
  10:00--10:50 & \textbf{1. Foundations:} potential outcomes and the fundamental problem & Lecture \\
  10:50--11:00 & \emph{Break} & \\
  11:00--11:30 & \textbf{2a. Experiments:} why randomization works & Lecture \\
  11:30--12:00 & \textbf{2b. Lab 1:} simulating potential outcomes and randomization & Lab \\
  \midrule
  12:00--13:00 & \emph{Lunch} & \\
  \midrule
  13:00--13:50 & \textbf{3a. Observational studies:} DAGs and designs that mimic experiments & Lecture \\
  13:50--14:20 & \textbf{3b. Lab 2:} difference-in-differences & Lab \\
  14:20--14:30 & \emph{Break} & \\
  14:30--14:55 & \textbf{4. Advanced topics} (overview) & Lecture \\
  14:55--15:00 & Wrap-up and where to go next & \\
  \bottomrule
\end{tabularx}

\smallskip
\noindent{\small\color{midgrey}Times are indicative and may be adjusted on the day.}

% ================= BLOCK 1 =================
\section{Block 1: Foundations}
\blockmeta{50 minutes}{Lecture and discussion}

\begin{itemize}
  \item \textbf{Correlation vs.\ causation:} what makes a research question \emph{causal}?
  \item \textbf{Counterfactual thinking:} what would have happened to \emph{this} unit otherwise?
  \item \textbf{Potential outcomes:} $Y_i(1)$, $Y_i(0)$, and the individual treatment effect $\tau_i = Y_i(1) - Y_i(0)$.
  \item \textbf{The fundamental problem of causal inference:} we only ever observe one potential outcome per unit, so causal inference is a missing-data problem.
  \item \textbf{Estimands:} the average treatment effect (ATE), the average treatment effect on the treated (ATT), and on the untreated (ATU).
  \item \textbf{Selection bias:} a naive difference in means equals the ATT plus selection bias,
    \[
      \underbrace{E[Y_i \mid D_i{=}1] - E[Y_i \mid D_i{=}0]}_{\text{naive difference}}
      = \underbrace{E[Y_i(1) - Y_i(0) \mid D_i{=}1]}_{\text{ATT}}
      + \underbrace{E[Y_i(0) \mid D_i{=}1] - E[Y_i(0) \mid D_i{=}0]}_{\text{selection bias}} .
    \]
  \item \textbf{The experimental ideal:} if we could design the perfect study, what would it look like?
\end{itemize}

\begin{reading}
\textbf{\faBookOpen\ Reading.}
\TheEffect, Ch.~1 (Designing Research), Ch.~5 (Identification), Ch.~10 (Treatment Effects).\\
\emph{Supplementary:} Cunningham, S. \emph{Causal Inference: The Mixtape}, Ch.~4 (Potential Outcomes Causal Model). Free online at \url{https://mixtape.scunning.com}.
\end{reading}

% ================= BLOCK 2 =================
\section{Block 2: Experiments}
\blockmeta{60 minutes}{30-minute lecture plus 30-minute lab}

\subsection{Lecture}
\begin{itemize}
  \item \textbf{Why randomization works:} treatment assignment becomes independent of the potential outcomes, so treatment and control groups are balanced on both observed and unobserved characteristics.
  \item \textbf{Estimation:} the difference in means; regression with covariates (to improve precision, not to identify the effect).
  \item \textbf{Experiments in the social sciences:} lab, field, and survey experiments (vignettes, conjoints).
  \item \textbf{Threats to experiments:} noncompliance (intention-to-treat vs.\ complier effects), attrition, spillovers and SUTVA, external validity.
  \item \textbf{Worked example:} Gerber, Green \& Larimer (2008), social pressure and voter turnout.
\end{itemize}

\begin{labbox}{Lab 1: Playing God with potential outcomes}
\begin{enumerate}[leftmargin=1.4em]
  \item Simulate a dataset in which we observe \emph{both} $Y_i(0)$ and $Y_i(1)$, and compute the true ATE.
  \item Hide one potential outcome per unit to make the fundamental problem concrete.
  \item Let units self-select into treatment and see how the naive estimate is biased.
  \item Randomize treatment instead. Repeat 1{,}000 times and plot the sampling distribution of the estimate.
  \item Check covariate balance, and add a covariate to a regression to see the gain in precision.
  \item \emph{Extension:} add noncompliance and compare the as-treated, intention-to-treat and complier average causal effect (CACE) estimates with the true effect among compliers.
\end{enumerate}
\end{labbox}

\begin{reading}
\textbf{\faBookOpen\ Reading.}
\TheEffect, Ch.~9 (Finding Front Doors), especially \S9.2 on randomized experiments.\\
\emph{Supplementary:} Gerber, A. S. \& Green, D. P. (2012). \emph{Field Experiments: Design, Analysis, and Interpretation}, Ch.~1--2.
\end{reading}

% ================= BLOCK 3 =================
\section{Block 3: Observational Studies}
\blockmeta{80 minutes}{50-minute lecture plus 30-minute lab}

\subsection{Lecture}
\begin{itemize}
  \item \textbf{When experiments aren't feasible:} ethical, practical and financial limits on randomization.
  \item \textbf{Directed acyclic graphs (DAGs):} nodes, edges, paths and back doors; confounders vs.\ mediators vs.\ colliders; selection bias as conditioning on a collider.
  \item \textbf{Designs that approximate the experimental ideal.} Each design is covered as: the intuition, the key assumption, one example, and what could break it.
  \begin{itemize}
    \item \emph{Difference-in-differences:} parallel trends.
    \item \emph{Regression discontinuity:} continuity of potential outcomes at the cutoff.
    \item \emph{Instrumental variables:} relevance, independence and the exclusion restriction (plus monotonicity for the LATE).
    \item \emph{Briefly:} regression and matching as ``selection on observables''.
  \end{itemize}
\end{itemize}

\begin{labbox}{Lab 2: Difference-in-differences}
\begin{itemize}[leftmargin=1.4em]
  \item \emph{Warm-up (optional):} draw a DAG with \texttt{dagitty} and find a valid adjustment set.
  \item \textbf{Application:} did New Jersey's 1992 minimum wage increase reduce employment? Using Card and Krueger's (1994) survey of fast-food restaurants in New Jersey and Pennsylvania, we build the treatment and post-period indicators, check that wages rose, estimate a $2\times2$ difference-in-differences model with \texttt{fixest}, and use high-wage New Jersey restaurants as a second comparison group. A simulation shows what an event study adds when there are more time periods.
\end{itemize}
\end{labbox}

\begin{reading}
\textbf{\faBookOpen\ Reading.}
\TheEffect, Ch.~6--8 (Causal Diagrams; Drawing Causal Diagrams; Causal Paths and Closing Back Doors) and Ch.~18 (Difference-in-Differences).
Skim Ch.~19 (Instrumental Variables) and Ch.~20 (Regression Discontinuity).
\end{reading}

% ================= BLOCK 4 =================
\section{Block 4: Advanced Topics}
\blockmeta{25 minutes}{Lecture overview}

A short tour of where the field has gone. For each topic: what problem it solves, the key idea, and one paper to start with.

\begin{itemize}
  \item \textbf{Staggered adoption and the problem with two-way fixed effects:} Goodman-Bacon (2021); Callaway \& Sant'Anna (2021).
  \item \textbf{Synthetic control:} Abadie, Diamond \& Hainmueller (2010).
  \item \textbf{Sensitivity analysis:} how strong would an unobserved confounder have to be to overturn a result? Cinelli \& Hazlett (2020), \texttt{sensemakr}.
  \item \textbf{Heterogeneous effects and machine learning:} causal forests (Wager \& Athey 2018); double/debiased machine learning (Chernozhukov et al.\ 2018).
  \item \textbf{Interference and mediation:} when SUTVA fails, and how to ask \emph{how} a treatment works. Imai, Keele, Tingley \& Yamamoto (2011).
\end{itemize}

\begin{reading}
\textbf{\faBookOpen\ Reading.}
\TheEffect, Ch.~21 (Partial Identification), Ch.~22 (A Gallery of Rogues: Other Methods), Ch.~23 (Under the Rug).
\end{reading}

% ================= PRACTICAL =================
\section{Practical Information}

\subsection{Software}
Labs use \textbf{R}. Before the workshop, please install R and RStudio and the following packages:
\begin{center}
  \texttt{tidyverse}\quad \texttt{fixest}\quad \texttt{panelView}\quad \texttt{dagitty}\quad \texttt{ggdag}
\end{center}
All code and data are in the course GitHub repository.

\subsection{Further resources}
\begin{itemize}
  \item Huntington-Klein, N. \emph{The Effect}. \url{https://theeffectbook.net}
  \item Cunningham, S. \emph{Causal Inference: The Mixtape}. \url{https://mixtape.scunning.com}
  \item Hern\'an, M. A. \& Robins, J. M. \emph{Causal Inference: What If}. \url{https://miguelhernan.org/whatifbook}
  \item Angrist, J. D. \& Pischke, J.-S. (2009). \emph{Mostly Harmless Econometrics}. Princeton University Press.
\end{itemize}

% ================= REFERENCES =================
\section{References}
{\small
\begin{itemize}[label={}, leftmargin=1.5em, itemindent=-1.5em]
  \item Abadie, A., Diamond, A., \& Hainmueller, J. (2010). Synthetic control methods for comparative case studies. \emph{Journal of the American Statistical Association}, 105(490), 493--505.
  \item Callaway, B., \& Sant'Anna, P. H. C. (2021). Difference-in-differences with multiple time periods. \emph{Journal of Econometrics}, 225(2), 200--230.
  \item Chernozhukov, V., et al.\ (2018). Double/debiased machine learning for treatment and structural parameters. \emph{The Econometrics Journal}, 21(1), C1--C68.
  \item Cinelli, C., \& Hazlett, C. (2020). Making sense of sensitivity: Extending omitted variable bias. \emph{Journal of the Royal Statistical Society: Series B}, 82(1), 39--67.
  \item Gerber, A. S., Green, D. P., \& Larimer, C. W. (2008). Social pressure and voter turnout. \emph{American Political Science Review}, 102(1), 33--48.
  \item Goodman-Bacon, A. (2021). Difference-in-differences with variation in treatment timing. \emph{Journal of Econometrics}, 225(2), 254--277.
  \item Imai, K., Keele, L., Tingley, D., \& Yamamoto, T. (2011). Unpacking the black box of causality: Learning about causal mechanisms from experimental and observational studies. \emph{American Political Science Review}, 105(4), 765--789.
  \item Wager, S., \& Athey, S. (2018). Estimation and inference of heterogeneous treatment effects using random forests. \emph{Journal of the American Statistical Association}, 113(523), 1228--1242.
\end{itemize}
}

\end{document}
